% ============================================================================
% 北京市建设工程质量第三检测所 — 检测报告 LaTeX 模板
% 依据《上传检测报告格式要求》(2016-01-01) 编写
% 编译: xelatex template.tex   (需 ctex / 系统含 SimSun / SimHei / Times New Roman)
% ============================================================================
\documentclass[UTF8, zihao=-4, a4paper]{ctexart}
% zihao=-4 即「小四」，对应正文字号；图表内字号在局部用 \zihao{5}/\zihao{-5}

% ---- 页面 ----
\usepackage[top=2.5cm, bottom=2.5cm, left=3.0cm, right=2.0cm,
            headheight=14pt, headsep=2cm, footskip=1.5cm]{geometry}

% ---- 字体（汉字 宋体；强调 黑体；西文 Times New Roman） ----
\setCJKmainfont{SimSun}[BoldFont=SimHei, ItalicFont=KaiTi]
\setCJKsansfont{SimHei}
\setmainfont{Times New Roman}

% ---- 行距与缩进 ----
\linespread{1.5}
\setlength{\parindent}{2em}            % 首行缩进 2 个汉字字符
\setlength{\parskip}{0pt}              % 段间距 0

% ---- 标题（一/二/三 级） ----
\usepackage{titlesec}
% 一级：对中、小三号 宋体加粗、段前 2 行段后 1 行、单倍行间距
\titleformat{\section}
  {\centering\zihao{-3}\bfseries\linespread{1.0}\selectfont}
  {\thesection}{1em}{}
\titlespacing*{\section}{0pt}{2\baselineskip}{1\baselineskip}
% 二级：左对齐、四号 黑体、段前段后各 1 行、单倍行间距
\titleformat{\subsection}
  {\zihao{4}\heiti\linespread{1.0}\selectfont}
  {\thesubsection}{1em}{}
\titlespacing*{\subsection}{0pt}{1\baselineskip}{1\baselineskip}
% 三级：左对齐、小四号 宋体加粗、段前 1 行段后 0.5 行、单倍行间距
\titleformat{\subsubsection}
  {\zihao{-4}\bfseries\linespread{1.0}\selectfont}
  {\thesubsubsection}{1em}{}
\titlespacing*{\subsubsection}{0pt}{1\baselineskip}{0.5\baselineskip}

% ---- 图表标题（五号 宋体、不加粗不倾斜、对中） ----
\usepackage{caption}
\captionsetup{labelsep=quad, justification=centering,
              labelfont={}, textfont={}, font={stretch=1.0}}

% ---- 封面绝对定位用 ----
\usepackage{tikz}
\usetikzlibrary{calc}

% ---- 图表节内编号（表1-1、表1.1-1 形式由用户在 \caption 里手动写出） ----
\usepackage{graphicx}
\usepackage{array, booktabs, longtable, makecell, multirow}

% ---- 页眉页脚（宋体 小三号） ----
\usepackage{fancyhdr}
\pagestyle{fancy}
\fancyhf{}
\renewcommand{\headrulewidth}{0.4pt}
\renewcommand{\footrulewidth}{0pt}
\fancyhead[C]{\zihao{-3}\songti 北京市建设工程质量第三检测所有限责任公司\\
              工程质量检测报告}
\fancyfoot[L]{\zihao{-3}\songti 报告编号：\reportno}
\fancyfoot[R]{\zihao{-3}\songti 第~\thepage~页\,/\,共~\pageref{LastPage}~页}
\usepackage{lastpage}

% 封面专用页面样式（无页眉页脚）
\fancypagestyle{cover}{\fancyhf{}\renewcommand{\headrulewidth}{0pt}}

% ---- 超链接（隐藏框） ----
\usepackage[hidelinks]{hyperref}

% ---- 实用宏 ----
\newcommand{\reportno}{京建质检J3—G字2016第（****）号}
\newcommand{\projname}{*******工程}
\newcommand{\client}{******公司}
\newcommand{\insttype}{一般委托}
\newcommand{\issuer}{北京市建设工程质量第三检测所有限责任公司}
\newcommand{\reportdate}{2015年01月01日}

% ============================================================================
\begin{document}

% ============================================================================
% 封面
% ============================================================================
\thispagestyle{cover}
% 标题"检 测 报 告"顶端距页面上边缘 60mm，初号 黑体，字间空 1 格 —— 用 tikz 绝对定位
\begin{tikzpicture}[remember picture, overlay]
  \node[anchor=north, inner sep=0pt]
       at ([yshift=-60mm]current page.north)
       {\zihao{0}\heiti 检\quad 测\quad 报\quad 告};
\end{tikzpicture}%
\vspace*{55mm}
\begin{center}
  {\zihao{2}\songti \reportno}\\[3cm]

  {\zihao{3}\songti
  \begin{tabular}{rl}
    工\ 程\ 名\ 称： & \projname \\[12pt]
    委\ 托\ 单\ 位： & \client  \\[12pt]
    检\ 测\ 类\ 别： & \insttype \\
  \end{tabular}}\\[3cm]

  {\zihao{2}\songti \issuer}\\[1cm]
  {\zihao{3}\songti \reportdate}
\end{center}
\clearpage

% ============================================================================
% 签字页（页码从此开始计为第 1 页）
% ============================================================================
\setcounter{page}{1}
\thispagestyle{fancy}

\begin{center}{\zihao{-2}\bfseries\songti 检测报告基本信息}\end{center}
\vspace{1em}
\renewcommand{\arraystretch}{1.6}
\begin{longtable}{|p{3cm}|p{11cm}|}
\hline 工程名称 & \projname \\\hline
工程地点 & 北京市海淀区 \\\hline
检测项目 & 混凝土抗压强度 \\\hline
施工单位 & ****** \\\hline
设计单位 & ****** \\\hline
监理单位 & ****** \\\hline
检测原因 & 委托检测 \\\hline
抽样依据 & 委托方指定 \\\hline
方法依据 & 《回弹法、超声回弹综合法检测泵送混凝土强度技术规程》DBJ/T 01-78-2003 \\\hline
判定依据 & 《回弹法、超声回弹综合法检测泵送混凝土强度技术规程》DBJ/T 01-78-2003 \\\hline
检测仪器 & ZBL-S220+S30901013-0867 数显混凝土回弹仪；LG1-19-08 碳化深度仪 \\\hline
检测时间 & 2015年01月01日 \\\hline
检测环境 & 晴~3\,$^{\circ}$C \\\hline
检测人员 & 胡运、李瑾、张恒源 \\\hline
检测结论 & 本次检测的 \projname 一至二层梁、板、柱 30 个构件现龄期混凝土抗压强度推定值达到设计强度标准。 \\\hline
\end{longtable}

\vspace{1em}
\begin{tabular}{p{3cm}p{4cm}p{3cm}p{4cm}}
批\ 准\ 人： & \makebox[3cm]{\hrulefill} & 审\ 核\ 人： & \makebox[3cm]{\hrulefill} \\[18pt]
主\ 检\ 人： & \makebox[3cm]{\hrulefill} & 报告日期： & \reportdate \\
\end{tabular}

\vspace{2em}
\begin{flushright}\zihao{-4}\issuer\end{flushright}
\clearpage

% ============================================================================
% 目录（页码 ≥ 15 页时启用）
% ============================================================================
% \tableofcontents
% \clearpage

% ============================================================================
% 正文 —— 以 二、回弹法检测混凝土强度报告 为示例
% （切换报告类型只需替换本节内容；通用样式不变）
% ============================================================================

\section{概况}

\subsection{工程概况}
\projname 位于北京市海淀区。建设单位为 \client 。综合楼建筑面积 4747.5\,m$^2$，
为钢筋混凝土框架结构，地下一层，地上三层，建筑高度 14.1\,m。\ldots\ldots
受 \client 委托，\issuer 对该工程联合检修库一至二层梁、板、柱 30 个构件混凝土抗压强度
进行了检测。

\subsection{检测及判定依据}
《回弹法、超声回弹综合法检测泵送混凝土强度技术规程》DBJ/T 01-78-2003。

\subsection{检测人员}
胡运、李瑾、张恒源。

\subsection{检测仪器}
\begin{table}[h]\centering
\caption*{\zihao{5}\songti 表1.4\quad 检测仪器}
\zihao{5}
\begin{tabular}{|l|c|c|}
\hline
检测仪器名称   & 数显混凝土回弹仪          & 碳化深度仪          \\\hline
检测仪器编号   & ZBL-S220+S30901013-0867 & LG1-19-08          \\\hline
检定证书编号   & 京西字第14037414号       & 京西字第14020308号  \\\hline
仪器使用有效期 & 2014.8.5\,$\sim$\,2015.2.4 & 2014.5.13\,$\sim$\,2015.5.12 \\\hline
检测精度       & 2                       & 0.25\,mm           \\\hline
\end{tabular}
\end{table}

\subsection{检测环境}
晴~3\,$^{\circ}$C。

\section{混凝土抗压强度检测}

\subsection{检测方法}
按照《回弹法、超声回弹综合法检测泵送混凝土强度技术规程》DBJ/T 01-78-2003 要求，
在每一个待测构件侧面或底面均匀布置 10 个测区，每个测区大小约为 0.04\,m$^2$，
在每个测区内测取 16 个回弹测点，并选取 30\% 的测区测量混凝土碳化深度值。

\subsection{抽样依据}
检测数量由委托方指定，现场对该工程联合检修库一至二层梁、板、柱共 30 个构件
进行混凝土抗压强度检测。

\subsection{判定依据}
依据《回弹法、超声回弹综合法检测泵送混凝土强度技术规程》DBJ/T 01-78-2003
进行单个构件混凝土抗压强度推定。

\subsection{检测结果}
本次检测的 \projname 一至二层梁、板、柱 30 个构件现龄期混凝土抗压强度推定值
达到设计强度标准。混凝土抗压强度检测结果详见表 2.4，回弹测区位置示意图见图 3-1\,$\sim$\,图 3-3。

\begin{table}[h]\centering
\caption*{\zihao{5}\songti 表2.4\quad 混凝土抗压强度检测结果}
\zihao{5}
\renewcommand{\arraystretch}{1.4}
\begin{tabular}{|c|l|c|c|c|c|}
\hline
\makecell{序\\号} & 构件名称 & \makecell{测区强度\\换算值\\(MPa)} & \makecell{换算值\\平均值\\(MPa)} & \makecell{换算值\\标准差\\(MPa)} & \makecell{构件抗压\\强度推定值\\(MPa)} \\\hline
1 & 一层 2-9/2-A$\sim$2-B 梁  & ~ & ~ & ~ & ~ \\\hline
2 & 一层 2-11/2-A$\sim$2-B 梁 & ~ & ~ & ~ & ~ \\\hline
3 & 一层 2-14/2-A$\sim$2-B 梁 & ~ & ~ & ~ & ~ \\\hline
4 & 一层 2-16/2-A$\sim$2-B 梁 & ~ & ~ & ~ & ~ \\\hline
5 & 一层 2-17/2-A$\sim$2-B 梁 & ~ & ~ & ~ & ~ \\\hline
\end{tabular}
\end{table}

\section{附图}

\begin{figure}[h]\centering
% \includegraphics[width=0.7\linewidth]{fig/3-1.png}
\framebox[0.7\linewidth]{\rule{0pt}{4cm}~图3-1 占位~}
\caption*{\zihao{5}\songti 图3-1\quad 梁回弹测区位置示意图（单位：cm）}
\end{figure}

\begin{figure}[h]\centering
\framebox[0.7\linewidth]{\rule{0pt}{4cm}~图3-2 占位~}
\caption*{\zihao{5}\songti 图3-2\quad 板回弹测区位置示意图（单位：cm）}
\end{figure}

\begin{figure}[h]\centering
\framebox[0.7\linewidth]{\rule{0pt}{4cm}~图3-3 占位~}
\caption*{\zihao{5}\songti 图3-3\quad 柱回弹测区位置示意图（单位：cm）}
\end{figure}

\begin{figure}[h]\centering
\framebox[0.7\linewidth]{\rule{0pt}{4cm}~现场照片占位~}
\caption*{\zihao{5}\songti 图3-4\quad 回弹法检测混凝土强度现场检测照片}
\end{figure}

\begin{center}（本页以下空白）\end{center}

\end{document}
